\chapter{Мультимодальная реконструкция визуальных стимулов по сигналам фМРТ--ЭЭГ}\label{ch:bimodal}

В данной главе рассматривается третья постановка задачи декодирования, расширяющая построенное пространственно-временное описание на вторую модальность: по одновременным сигналам фМРТ и ЭЭГ реконструируются визуальные стимулы. Нейросетевая модель обучается на записях всех пациентов, а межсубъектные различия учитываются индивидуальными компонентами энкодеров. Совместные представления модальностей контрастивно выравниваются с CLIP-пространством, после чего реконструкция выполняется двухстадийной диффузионной генерацией. Результаты главы опубликованы в работах~\cite{dorin2026decoding, dorin2025mmro}.

Полный код реализации, использованный для получения результатов главы, доступен в открытом репозитории:
\begin{center}
    {\small \url{https://github.com/intsystems/CreationOfIntelligentSystems_Simultaneous_fMRI-EEG}}
\end{center}

\section{Постановка задачи}\label{sec2:bimodal}

Задан датасет триплетов (фМРТ, ЭЭГ, изображение):
\begin{equation*}\mathfrak{D} = \{(\mathbf{F}_i, \mathbf{E}_i, \mathbf{I}_i) \mid i = 1, \ldots, N \}, \quad (\mathbf{F}_i, \mathbf{E}_i) \in \mathcal{X}_f \times \mathcal{X}_e, \quad \mathbf{I}_i \in \mathcal{Y},
\end{equation*}
где $\mathbf{F}_i \in \mathcal{X}_f$~--- фМРТ-скан; $\mathbf{E}_i \in \mathcal{X}_e \subseteq \mathbb{R}^{K \times T}$~--- ЭЭГ-запись, где $K$~--- число каналов, $T$~--- длина временного окна; $\mathbf{I}_i \in \mathcal{Y} \subseteq \mathbb{R}^{H \times W \times C}$~--- изображение с высотой $H$, шириной $W$ и числом каналов $C$. Требуется построить отображение $\mathbf{f}$ из совместного пространства данных фМРТ--ЭЭГ в пространство изображений:
\begin{equation}\label{eq:target}
    \mathbf{f}:  \mathcal{X}_f \times \mathcal{X}_e \rightarrow \mathcal{Y}.
\end{equation}

\section{Метод реконструкции}\label{sec3:bimodal}
До появления недавних мультимодальных наборов данных с визуальными стимулами, таких как~\cite{Telesford2023}, обработка ЭЭГ и фМРТ, как правило, проводилась раздельно. Новые датасеты открывают возможность интегрировать оба типа сигналов в единую архитектуру. По аналогии с предшествующими работами, в настоящей работе используется двухэтапный метод декодирования мозговой активности: 1)~контрастивное выравнивание сигналов фМРТ и ЭЭГ с эмбеддингами изображений; 2)~обуславливание предобученной модели генерации изображений на уточненный априорной моделью эмбеддинг сигналов мозга.

На Рисунке~\ref{fig:scheme} представлена схема декодирования изображений. Процесс строится из двух основных этапов: контрастивное обучение комбинированных представлений фМРТ--ЭЭГ сигналов и реконструкция стимулов. Предварительно записи ЭЭГ приводятся к фиксированному набору каналов восстановлением отсутствующих электродов, описанным в разделе~\ref{sec3:recover}.

\subsection{Восстановление отсутствующих каналов ЭЭГ}\label{sec3:recover}
На этапе предобработки решается задача восстановления пропусков в записях ЭЭГ: сигналы незарегистрированных электродов восстанавливаются методами пространственной интерполяции на основе данных ближайших доступных каналов и известной геометрии расположения датчиков.

Полный набор электродов системы регистрации включает $K$ каналов. Каждому каналу $\{c_k\}_{k=1}^K$ соответствуют фиксированные пространственные координаты $\mathbf{p}_k \in \bbR^3$ на поверхности скальпа. Множество фактически зарегистрированных каналов в $i$-й записи обозначим как $\mathcal{C}_i \subseteq \{c_1, \ldots, c_K\},~ i=1,\ldots, N$. От записи к записи набор отсутствующих каналов меняется, поэтому ограничиваться лишь общим для всех записей подмножеством каналов нельзя~--- это привело бы к существенной потере информации. При этом энкодер ЭЭГ $\mathcal{E}_e$ принимает на вход тензор фиксированной размерности, без чего невозможна обработка обучающей выборки батчами. Отсюда и вытекает необходимость восстановить сигналы недостающих каналов до полного набора.

Для упрощения нотации далее рассматривается единичная запись ЭЭГ $\mathbf{E}$, при этом индекс $i$ объекта выборки опускается. Обозначим через $\mathcal{C} \subset \{c_1, \dots, c_K\}$~--- подмножество каналов, для которых доступны временные ряды $\mathbf{e}_k \in \bbR^T$. Таким образом, наблюдаемая часть записи задается как $\mathbf{E} = \{\mathbf{e}_k\}_{k \in \mathcal{C}}$.
Для унификации формата данных каждому отсутствующему каналу $c_k \notin \mathcal{C}$ назначается восстановленный сигнал $\widehat{\mathbf{e}}_k \in \bbR^T$, рассчитываемый как взвешенное среднее значений $n$ пространственно ближайших зарегистрированных электродов:
\begin{equation}\label{eq:recover}
\widehat{\mathbf{e}}_k = \sum_{j \in \mathcal{S}_k} w_{kj}\, \mathbf{e}_j, \qquad \sum_{j \in \mathcal{S}_k} w_{kj} = 1,
\end{equation}
где $\mathbf{e}_j \in \bbR^T$~--- временной ряд канала $c_j$; $\mathcal{S}_k$~--- множество индексов $n$ ближайших к $c_k$ доступных электродов, определяемое в соответствии с выбранной пространственной метрикой; $w_{kj} \geq 0$~--- неотрицательные нормированные веса.
В исследовании рассматриваются четыре стратегии определения весов $\{w_{kj}\}$ и формирования множества соседей $\mathcal{S}_k$:
\begin{itemize}
    \item \textbf{Zero}, нулевое заполнение: базовый метод, при котором отсутствующий сигнал заменяется нулевым вектором:
    \begin{equation*}
    \widehat{\mathbf{e}}_k = \mathbf{0} \in \bbR^T.
    \end{equation*}
    \item \textbf{NN}, метод одного ближайшего соседа (англ. Nearest Neighbor, NN): сигнал отсутствующего канала заменяется сигналом ближайшего зарегистрированного электрода:
    \begin{equation*}
    \widehat{\mathbf{e}}_k = \mathbf{e}_{j^*}, \qquad j^* = \argmin_{j: c_j \in \mathcal{C}} \rho_{\text{eucl}}(c_k, c_j), \quad \rho_{\text{eucl}}(c_k, c_j) = \|\mathbf{p}_k - \mathbf{p}_j\|_2.
    \end{equation*}
    что соответствует $n = 1$ и $w_{kj^*} = 1$ в~\eqref{eq:recover}.
    \item \textbf{KNN-Euclidean}, метод ближайших соседей (англ. $k$-Nearest Neighbors, KNN): взвешенное усреднение по $n$ ближайшим соседям в евклидовой метрике; веса определяются как величины, обратно пропорциональные расстоянию между электродами:
    \begin{equation*}
    w_{kj} = \frac{1 / \rho_{\text{eucl}}(c_k, c_j)}{\displaystyle\sum_{l \in \mathcal{S}_k} 1 / \rho_{\text{eucl}}(c_k, c_l)}, \qquad \rho_{\text{eucl}}(c_k, c_j) = \|\mathbf{p}_k - \mathbf{p}_j\|_2.
    \end{equation*}
    \item \textbf{KNN-Spherical}: то же взвешенное усреднение по $n$ ближайшим соседям, но в сферической метрике~--- евклидово расстояние здесь заменено геодезическим вдоль поверхности скальпа, которая аппроксимируется сферой радиуса $R$:
    \begin{equation*}
    \rho_{\text{sph}}(c_k, c_j) = R \cdot \arccos\!\left( \frac{\langle \mathbf{p}_k, \mathbf{p}_j \rangle}{\|\mathbf{p}_k\|_2 \, \|\mathbf{p}_j\|_2} \right),
    \end{equation*}
    веса при этом полагаются $w_{kj} \propto 1 / \rho_{\text{sph}}(c_k, c_j)$ и затем нормируются: $\sum_{j \in \mathcal{S}_k} w_{kj} = 1$.
\end{itemize}
Выполнение процедуры~\eqref{eq:recover} для всех отсутствующих каналов приводит запись ЭЭГ к единому матричному виду $\mathbf{E} \in \bbR^{K \times T}$ с фиксированным числом каналов $K$. Тем самым размерность входного тензора согласуется с архитектурой энкодера $\mathcal{E}_e$, что и делает возможной последующую батчевую обработку.

\paragraph{Вероятностное обобщение: маскированное согласование потоков.}
Интерполяция~\eqref{eq:recover} дает точечную оценку пропущенного сигнала и использует только пространственную близость электродов. Ниже она достраивается до вероятностной модели, которая восстанавливает условное распределение пропусков при наблюдаемой части записи. Модель строится на согласовании потоков (англ. flow matching)~\cite{lipman2023flow}: непрерывный нормализующий поток~\cite{chen2018neural} переносит гауссовский шум в данные, а обучается регрессией на условное векторное поле, без решения дифференциальных уравнений на этапе обучения.

Введем бинарную маску наблюдаемости $\mathbf{M} \in \{0,1\}^{K \times T}$, где $M_{k,t} = 1$ для зарегистрированных значений, и ее дополнение $\bar{\mathbf{M}} = \mathbf{1} - \mathbf{M}$; при выпадении канала $c_k$ нулевой является вся строка $k$. Наблюдаемая и пропущенная части записи равны $\mathbf{E}_{\mathrm{obs}} = \mathbf{M} \odot \mathbf{E}$ и $\mathbf{E}_{\mathrm{miss}} = \bar{\mathbf{M}} \odot \mathbf{E}$, где $\odot$~--- произведение Адамара. Геометрия монтажа задана матрицей расстояний $\mathbf{D} \in \bbR^{K \times K}$, $D_{kj} = \rho(c_k, c_j)$. Требуется по $(\mathbf{E}_{\mathrm{obs}}, \mathbf{M}, \mathbf{D})$ восстановить условное распределение $p(\mathbf{E}_{\mathrm{miss}} \mid \mathbf{E}_{\mathrm{obs}}, \mathbf{M})$. Обучающая выборка состоит из окон записей; маски при обучении семплируются независимо от сигнала из распределения, воспроизводящего типичные сценарии: выпадение целых каналов и короткие артефактные блоки~\cite{tashiro2021csdi}.

\textit{Маскированный путь и функция потерь.} Пусть $\boldsymbol{\varepsilon} \sim \mathcal{N}(\mathbf{0}, \mathbb{I})$~--- гауссовский шум той же размерности, что и запись, не зависящий от $(\mathbf{E}, \mathbf{M})$. Состояние в генеративное время $t \in [0,1]$ определяется покоординатно: наблюдаемые значения неподвижны, а пропущенные движутся по отрезку от шума к данным,
\begin{equation}\label{eq:imp_path}
\mathbf{E}_t = \mathbf{M} \odot \mathbf{E} + \bar{\mathbf{M}} \odot \big( (1-t)\,\boldsymbol{\varepsilon} + t\,\mathbf{E} \big).
\end{equation}
Производная состояния по $t$ равна $\bar{\mathbf{M}} \odot (\mathbf{E} - \boldsymbol{\varepsilon})$ и обращается в нуль на наблюдаемых позициях. Векторное поле $v_{\theta}(\mathbf{E}_t, t, \mathbf{E}_{\mathrm{obs}}, \mathbf{M}; \mathbf{D})$ задается нейросетью и обучается регрессией на это приращение по пропущенным позициям:
\begin{equation}\label{eq:imp_loss}
\mathcal{L}_{\mathrm{imp}}(\theta) = \mathbb{E}_{t, \mathbf{E}, \mathbf{M}, \boldsymbol{\varepsilon}} \Big\| \bar{\mathbf{M}} \odot \big( v_{\theta}(\mathbf{E}_t, t, \mathbf{E}_{\mathrm{obs}}, \mathbf{M}; \mathbf{D}) - (\mathbf{E} - \boldsymbol{\varepsilon}) \big) \Big\|_F^2,
\qquad t \sim \mathcal{U}[0,1].
\end{equation}
Наблюдаемая часть записи входит в модель дважды: как неподвижные координаты состояния и как условие во входах сети. Сеть состоит из блоков, чередующих свертку вдоль времени в каждом канале и внимание по каналам; в логиты внимания расстояния между электродами входят аддитивным смещением $B_{kj} = f_{\phi}(D_{kj})$, где $f_{\phi}$~--- небольшой перцептрон, общий для всех пар электродов, что дает модели индуктивное смещение в пользу физически близких электродов.

\textit{Генерация.} Восстановление записи с наблюдаемой частью $\mathbf{E}_{\mathrm{obs}}^{*}$ и маской $\mathbf{M}^{*}$ есть численное интегрирование обыкновенного дифференциального уравнения
\begin{equation}\label{eq:imp_ode}
\frac{d\mathbf{E}_t}{dt} = \bar{\mathbf{M}}^{*} \odot v_{\theta}\big(\mathbf{E}_t, t, \mathbf{E}_{\mathrm{obs}}^{*}, \mathbf{M}^{*}; \mathbf{D}\big), \qquad \mathbf{E}_0 = \mathbf{E}_{\mathrm{obs}}^{*} + \bar{\mathbf{M}}^{*} \odot \boldsymbol{\varepsilon}, \qquad t \colon 0 \to 1,
\end{equation}
методом Эйлера. Умножение поля на $\bar{\mathbf{M}}^{*}$ гарантирует, что наблюдаемые значения не меняются, поэтому состояния генерации лежат в том же множестве, что и обучающие состояния~\eqref{eq:imp_path}. Конечное состояние $\mathbf{E}_1$ есть восстановленная запись; повторение с независимыми реализациями $\boldsymbol{\varepsilon}$ дает ансамбль восстановлений, покоординатная медиана которого служит точечной оценкой, а разброс~--- оценкой неопределенности восстановления.

\textit{Корректность.} Следующая теорема показывает, что при точном векторном поле процедура~\eqref{eq:imp_ode} восстанавливает именно условное распределение пропусков. Зафиксируем $(\mathbf{E}_{\mathrm{obs}}, \mathbf{M})$ и обозначим через $\boldsymbol{\xi}_1 \in \bbR^{m}$ вектор пропущенных значений, где $m$~--- их число, с распределением $p(\cdot \mid \mathbf{E}_{\mathrm{obs}}, \mathbf{M})$, а через $\boldsymbol{\xi}_0 \sim \mathcal{N}(\mathbf{0}, \mathbb{I}_m)$~--- не зависящий от него вектор шума на тех же позициях. Положим $\boldsymbol{\xi}_t = (1-t)\,\boldsymbol{\xi}_0 + t\,\boldsymbol{\xi}_1$ и
\begin{equation}\label{eq:imp_field}
v^{*}(\boldsymbol{\xi}, t) = \mathbb{E}\big[ \boldsymbol{\xi}_1 - \boldsymbol{\xi}_0 \mid \boldsymbol{\xi}_t = \boldsymbol{\xi} \big], \qquad \boldsymbol{\xi} \in \bbR^{m}, \quad t \in [0,1].
\end{equation}

\begin{theorem}\label{thm:imp_correct}
Пусть $\mathbb{E}\|\boldsymbol{\xi}_1\|_2 < \infty$, а поле~\eqref{eq:imp_field} непрерывно по $(\boldsymbol{\xi}, t)$ и липшицево по $\boldsymbol{\xi}$ равномерно по $t$. Тогда
\begin{enumerate}
    \item[(а)] поле $v^{*}$ есть минимизатор функции потерь~\eqref{eq:imp_loss} по всем измеримым функциям состояния, времени и условия $(\mathbf{E}_{\mathrm{obs}}, \mathbf{M})$;
    \item[(б)] решение уравнения $\dot{\boldsymbol{\xi}}(t) = v^{*}(\boldsymbol{\xi}(t), t)$ с начальным условием $\boldsymbol{\xi}(0) = \boldsymbol{\xi}_0$ существует, единственно и $\boldsymbol{\xi}(1) \sim p(\cdot \mid \mathbf{E}_{\mathrm{obs}}, \mathbf{M})$; это уравнение есть ограничение~\eqref{eq:imp_ode} на пропущенные позиции при $v_{\theta} = v^{*}$.
\end{enumerate}
\end{theorem}
\begin{proof}
(а) При фиксированных условии и времени $t$ слагаемые функции потерь~\eqref{eq:imp_loss} на наблюдаемых позициях обнулены маской $\bar{\mathbf{M}}$, а пропущенные позиции состояния~\eqref{eq:imp_path} образуют вектор $\boldsymbol{\xi}_t$. Поэтому функция потерь есть среднеквадратичная ошибка предсказания вектора $\boldsymbol{\xi}_1 - \boldsymbol{\xi}_0$ по наблюдению $\boldsymbol{\xi}_t$. Среди всех измеримых функций наблюдения среднеквадратичную ошибку минимизирует условное математическое ожидание~\eqref{eq:imp_field}, единственное с точностью до множества нулевой вероятности; усреднение по $t$ и по условию не меняет минимизатора, поскольку оба подаются сети на вход.

(б) Процесс $\boldsymbol{\xi}_t$ линеен по $t$: $\dot{\boldsymbol{\xi}}_t = \boldsymbol{\xi}_1 - \boldsymbol{\xi}_0$, причем $\boldsymbol{\xi}_0 \sim \mathcal{N}(\mathbf{0}, \mathbb{I}_m)$ и $\boldsymbol{\xi}_1 \sim p(\cdot \mid \mathbf{E}_{\mathrm{obs}}, \mathbf{M})$. Пусть $q_t$~--- распределение $\boldsymbol{\xi}_t$. Для любой гладкой финитной функции $\varphi$
\begin{equation*}
\frac{d}{dt}\, \mathbb{E}\,\varphi(\boldsymbol{\xi}_t) = \mathbb{E}\big\langle \nabla\varphi(\boldsymbol{\xi}_t),\, \boldsymbol{\xi}_1 - \boldsymbol{\xi}_0 \big\rangle = \mathbb{E}\big\langle \nabla\varphi(\boldsymbol{\xi}_t),\, v^{*}(\boldsymbol{\xi}_t, t) \big\rangle,
\end{equation*}
где второе равенство следует из определения условного ожидания~\eqref{eq:imp_field}, а дифференцирование под знаком ожидания законно в силу $\mathbb{E}\|\boldsymbol{\xi}_1 - \boldsymbol{\xi}_0\|_2 < \infty$ и ограниченности $\nabla\varphi$. Значит, семейство $(q_t)_{t \in [0,1]}$ является слабым решением уравнения непрерывности $\partial_t q_t + \nabla \cdot (q_t\, v^{*}(\cdot, t)) = 0$ с начальным условием $q_0 = \mathcal{N}(\mathbf{0}, \mathbb{I}_m)$. Поле $v^{*}$ липшицево по $\boldsymbol{\xi}$ и непрерывно по $t$, поэтому уравнение $\dot{\boldsymbol{\xi}} = v^{*}(\boldsymbol{\xi}, t)$ имеет единственное решение $\boldsymbol{\xi}(t) = \Phi_t(\boldsymbol{\xi}(0))$ при любом начальном условии, и образ $(\Phi_t)_{\#} q_0$ начального распределения также является слабым решением того же уравнения непрерывности с тем же начальным условием. Для липшицевых полей слабое решение уравнения непрерывности единственно в классе слабо непрерывных семейств вероятностных мер~\cite[гл.~8]{ambrosio2008gradient}, откуда $(\Phi_t)_{\#} q_0 = q_t$ при всех $t$. При $t = 1$ получаем $\boldsymbol{\xi}(1) = \Phi_1(\boldsymbol{\xi}_0) \sim q_1$, а $q_1$~--- распределение $\boldsymbol{\xi}_1$, то есть $p(\cdot \mid \mathbf{E}_{\mathrm{obs}}, \mathbf{M})$. Наконец, на наблюдаемых позициях правая часть~\eqref{eq:imp_ode} равна нулю, а на пропущенных при $v_{\theta} = v^{*}$ совпадает с $v^{*}$, что и дает последнее утверждение.
\end{proof}

В экспериментах настоящей главы записи ЭЭГ приводятся к полному набору каналов детерминированным восстановлением~\eqref{eq:recover}; вероятностная модель описана как его обобщение, и ее применение к записям ЭЭГ составляет предмет дальнейшей работы.

\subsection{Контрастивное обучение комбинированных представлений}\label{sec3:contrastive}
\begin{figure}[h]
    \centering
    \includegraphics[width=1.\textwidth]{figure_1.pdf}
    \caption{Схема декодирования изображений на основе сигналов фМРТ--ЭЭГ. Процесс декодирования реализуется в два этапа. На первом этапе проводится контрастное обучение энкодера сигналов фМРТ--ЭЭГ, в ходе которого представления мозговых сигналов выравниваются с замороженными эмбеддингами изображений, полученными с помощью модели CLIP. На втором этапе осуществляется двухстадийная генерация реконструкций: сначала обучается диффузионная априорная модель для уточнения эмбеддингов сигналов мозга, а затем предобученная диффузионная модель, обусловленная на CLIP-эмбеддинги, генерирует реконструкции изображений}
    \label{fig:scheme}
\end{figure}
\paragraph{Энкодер сигналов фМРТ--ЭЭГ.}
Архитектура энкодера состоит из трех модулей: энкодера фМРТ, энкодера ЭЭГ и модуля объединения. Для обеспечения персонализированного представления мозговой активности каждый модуль включает индивидуализирующие компоненты, позволяющие явно моделировать межсубъектные различия~--- критически важный аспект при работе с нейросигналами. Архитектуры энкодеров фМРТ и ЭЭГ изображены на Рисунке~\ref{fig:encoders}.

\paragraph{Энкодер ЭЭГ.}
Энкодер ЭЭГ через конкатенацию объединяет исходные многоканальные данные с индивидуальными обучаемыми эмбеддингами пациентов, получая единое векторное представление. Архитектура энкодера включает последовательность трех сверточных слоев: одномерную свертку по временному измерению, одномерную свертку по каналам, то есть по пространству, и свертку для сжатия в вектор. Полученный эмбеддинг пропускается через голову в виде многослойного перцептрона (англ. Multilayer Perceptron, MLP), формируя итоговое представление. Формально, для батча данных ЭЭГ $\mathbf{E} \in \mathbb{R}^{B \times K \times T}$ энкодер $\mathcal{E}_{e}$ генерирует эмбеддинги:
\begin{equation*}
\mathbf{Z}_{e} = \mathcal{E}_{e}(\mathbf{E}) \in \mathbb{R}^{B \times d_{e}},
\end{equation*}
где $d_{e}$~--- размерность латентного пространства ЭЭГ.

\paragraph{Предобработка фМРТ.}
Из всего объема $X \times Y \times Z$ вокселей лишь относительно небольшая часть несет информацию о реакции мозга на визуальный стимул: значительную долю занимают фоновые воксели, значения которых либо близки к нулю, либо стабильны во времени и не зависят от стимула. Прямая подача всего тензора в обучаемую модель приводит к расточительной параметризации и снижает отношение сигнал/шум. Поэтому на этапе предобработки строится индивидуальная для каждого пациента бинарная маска, выделяющая только информативные воксели.

Пусть $F_{xyz}(t)$~--- значение вокселя с координатами $(x, y, z)$ в момент $t$ временного ряда фМРТ длиной $\mathcal{T}$. Для каждого вокселя вычисляется суммарная абсолютная разность значений между соседними кадрами:
\begin{equation*}\label{eq:mask_score}
m_{xyz} = \sum_{t=2}^{\mathcal{T}} \left| F_{xyz}(t) - F_{xyz}(t-1) \right|.
\end{equation*}
Величина $m_{xyz}$ интерпретируется как накопленная во времени вариация сигнала в данном вокселе: для фоновых вокселей $m_{xyz}$ близка к нулю, для активно реагирующих на стимул~--- существенно положительна. К одномерному распределению величин $\{m_{xyz}\}$ применяется метод Оцу~\cite{otsu1975threshold}, автоматически подбирающий порог между фоновыми и информативными вокселями. Алгоритм перебирает все допустимые пороги $\theta$, разбивая при каждом из них воксели на два класса~--- $C_0(\theta) = \{(x, y, z) : m_{xyz} \le \theta\}$ и $C_1(\theta) = \{(x, y, z) : m_{xyz} > \theta\}$~--- с долями $\omega_0(\theta), \omega_1(\theta) = 1 - \omega_0(\theta)$ и средними $\mu_0(\theta), \mu_1(\theta)$. Межклассовая дисперсия имеет вид
\begin{equation*}\label{eq:otsu_var}
\sigma^2(\theta) = \omega_0(\theta)\, \omega_1(\theta)\, \big[\mu_0(\theta) - \mu_1(\theta)\big]^2,
\end{equation*}
оптимальным же считается порог, доставляющий ей максимум:
\begin{equation*}\label{eq:otsu_threshold}
\theta^* = \argmax_{\theta} \sigma^2(\theta).
\end{equation*}
Найденный порог задает бинарную маску
\begin{equation}\label{eq:mask}
\mathcal{M}_{xyz} = \mathds{1}\{m_{xyz} > \theta^*\}, \qquad \mathcal{M} \in \{0, 1\}^{X \times Y \times Z},
\end{equation}
свою для каждого пациента; число информативных вокселей $\sum_{xyz} \mathcal{M}_{xyz}$ определяется анатомией и индивидуальной динамикой BOLD-сигнала. После этого для каждого фМРТ-скана $\mathbf{F}$ воксели со значением $\mathcal{M}_{xyz} = 1$ собираются в плотный вектор.

\paragraph{Энкодер фМРТ.}
Индивидуальные особенности пациентов энкодер фМРТ учитывает с помощью персонализированной линейной проекции. Вектор информативных вокселей, отобранных по бинарной маске $\mathcal{M}$, пропускается через индивидуальный для каждого пациента линейный слой, проецирующий признаки в общее латентное пространство размерности 4096. Завершающим этапом является обработка эмбеддингов через MLP с остаточными соединениями. Формально, для батча данных фМРТ $\mathbf{F} \in \mathbb{R}^{B \times X \times Y \times Z}$ энкодер $\mathcal{E}_{f}$ преобразует исходный тензор в латентное представление:
\begin{equation*}
\mathbf{Z}_{f} = \mathcal{E}_{f}(\mathbf{F}) \in \mathbb{R}^{B \times d_{f}},
\end{equation*}
где $d_{f}$~--- размерность латентного пространства фМРТ.

\paragraph{Модуль объединения.}
Полученные эмбеддинги $\mathbf{Z}_e$ и $\mathbf{Z}_f$ подаются на вход обучаемого модуля объединения $\mathcal{F}$, который интегрирует высокочастотную информацию из ЭЭГ и пространственно детализированные данные фМРТ, формируя совместные эмбеддинги:
\begin{equation*}
  \mathbf{Z}_c = \mathcal{F}(\mathbf{Z}_f, \mathbf{Z}_e) \in \mathbb{R}^{B \times d}.
\end{equation*}
В данной работе модуль объединения реализован как конкатенация эмбеддингов с последующей обработкой через MLP для извлечения совместных представлений. В качестве альтернативной реализации $\mathcal{F}$ также рассматривалось скользящее усреднение эмбеддингов модальностей вида
\begin{equation*}\label{eq:ewma}
\mathbf{Z}_c = \alpha\, \mathbf{Z}_f + (1 - \alpha)\, \mathbf{Z}_e, \qquad \alpha \in [0, 1],
\end{equation*}
где $\alpha$~--- обучаемый скалярный параметр, балансирующий вклад фМРТ- и ЭЭГ-модальностей. Сравнительный анализ показал, что MLP более эффективно выделяет совместные признаки из пространств фМРТ--ЭЭГ, поэтому именно эта конфигурация принята в качестве основной.

\begin{figure}[h]
    \centering
    \includegraphics[width=0.9\textwidth]{figure_2.pdf}
    \caption{Архитектуры кодировщиков фМРТ и ЭЭГ. \textit{Слева:} энкодер ЭЭГ интегрирует индивидуальные обучаемые эмбеддинги пациентов с исходными многоканальными данными посредством конкатенации, формируя единое векторное представление; далее применяются одномерные свертки по временному и канальному измерениям, а также MLP-голова для генерации итоговых эмбеддингов. \textit{Справа:} энкодер фМРТ включает персонализированное маскирование мозга, векторизацию вокселей и проекцию в латентное пространство через индивидуальный линейный слой; завершающая обработка осуществляется с использованием MLP с остаточными соединениями}
    \label{fig:encoders}
\end{figure}

\paragraph{Эмбеддинги изображений.}
Для извлечения визуальных эмбеддингов используется предобученная модель CLIP ViT-H/14~\cite{radford2021learning}, которая генерирует эмбеддинги изображений $\mathbf{Z}_I \in \mathbb{R}^{B \times d}$. Данная модель обеспечивает высококачественные, семантически насыщенные и хорошо обобщающиеся эмбеддинги, а также совместима с множеством генеративных моделей, предобученных на CLIP-пространстве.

\paragraph{Контрастивное обучение энкодера сигналов мозга.}
Основной целью контрастивного обучения является сопоставление нейронных представлений фМРТ--ЭЭГ с визуальными эмбеддингами, извлеченными с помощью замороженного энкодера CLIP, как показано на Рисунке~\ref{fig:scheme}. Обучение происходит на триплетах (фМРТ, ЭЭГ, изображение) и включает оптимизацию трех ключевых компонентов: энкодера фМРТ $\mathcal{E}_f$, энкодера ЭЭГ $\mathcal{E}_e$ и модуля объединения $\mathcal{F}$. Для вычисления логитов близости используется косинусное расстояние между нормированными эмбеддингами:
\begin{equation*}\mathbf{Z}_I^n = \frac{\mathbf{Z}_I}{\|\mathbf{Z}_I\|}, \quad \mathbf{Z}_c^n = \frac{\mathbf{Z}_c}{\|\mathbf{Z}_c\|},
\end{equation*}
где $\mathbf{Z}_I$~--- эмбеддинги изображений, полученные с помощью замороженного CLIP, а $\mathbf{Z}_c$~--- объединенные эмбеддинги фМРТ--ЭЭГ. Матрица логитов $\mathbf{L}$ вычисляется как
\begin{equation*}\mathbf{L} = \mathbf{Z}_I^n \left(\mathbf{Z}_c^n\right)^{\top} \cdot e^{\tau},
\end{equation*}
где $\tau$~--- обучаемый параметр температуры. В процессе контрастивного обучения рассматриваются все возможные пары <<объединенный эмбеддинг мозга~--- эмбеддинг изображения>> для максимизации эффективности обучения. Важной особенностью подхода является наличие большого числа негативных примеров, что позволяет модели не только корректно сопоставлять положительные пары, но и четко различать негативные. Как показано на Рисунке~\ref{fig:contrastive_learning}, при размере батча $B$ формируется матрица логитов размером $B \times B$, где лишь $B$ пар, расположенных на главной диагонали, являются положительными соответствиями.

Функция потерь основывается на симметричной контрастивной оптимизации, аналогичной подходу CLIP~\cite{radford2021learning}, и представляет собой усредненную сумму двух кросс-энтропийных потерь:
\begin{equation}\label{eq:contrastive_loss}
    \mathcal{L}_{\text{contr}} = -\frac{1}{2B} \sum_{b=0}^{B-1} \left[ \log \left( \frac{\exp(\mathbf{L}_{bb})}{\sum_{j=0}^{B-1} \exp(\mathbf{L}_{bj})} \right) + \log \left( \frac{\exp(\mathbf{L}_{bb})}{\sum_{i=0}^{B-1} \exp(\mathbf{L}_{ib})} \right) \right],
\end{equation}
где $\mathbf{L}_{ij}$ обозначает элемент матрицы логитов $\mathbf{L}$ на позиции $(i, j)$.

\begin{figure}[h]
    \centering
    \includegraphics[width=0.85\textwidth]{figure_3.pdf}
    \caption{Иллюстрация контрастивного обучения и формирования батчей. Отрицательные примеры формируются автоматически. Косинусные расстояния между эмбеддингами замороженного CLIP и нейронными эмбеддингами используются в качестве логитов}
    \label{fig:contrastive_learning}
\end{figure}

\subsection{Процесс реконструкции стимулов}\label{sec3:gen}

Вводится совместное распределение изображений $\mathbf{x}$ и их скрытых визуальных представлений $\mathbf{z}_I$ при условии комбинированных нейронных эмбеддингов $\mathbf{z}_c$:
\begin{equation*}
p(\mathbf{x}, \mathbf{z}_I \mid \mathbf{z}_c) \propto p(\mathbf{z}_I \mid \mathbf{z}_c)\, p(\mathbf{x} \mid \mathbf{z}_I).
\end{equation*}
Здесь $p(\mathbf{z}_I \mid \mathbf{z}_c)$ моделирует согласование комбинированного представления фМРТ--ЭЭГ с высокоуровневым визуальным пространством CLIP на стадии~I, а $p(\mathbf{x} \mid \mathbf{z}_I)$ описывает генерацию изображения при условии уточненного визуального эмбеддинга на стадии~II.

\paragraph{Стадия I: диффузионная модель в пространстве CLIP.}
В пространстве визуальных эмбеддингов CLIP обучается легкая диффузионная модель вида
\begin{equation*}
\boldsymbol{\epsilon}_{\boldsymbol{\theta}}(\mathbf{z}_I^{(t)},\, t,\, \mathbf{z}_c),
\end{equation*}
где $\mathbf{z}_I^{(t)}$~--- зашумленная версия истинного эмбеддинга $\mathbf{z}_I$ на шаге $t$ прямого диффузионного процесса, а $\mathbf{z}_c$~--- комбинированный эмбеддинг, полученный модулем слияния. Обучение осуществляется по методу DDPM (англ. Denoising Diffusion Probabilistic Models, вероятностные диффузионные модели расшумления)~\cite{ho2020denoisingdiffusionprobabilisticmodels} с функцией потерь
\begin{equation}\label{eq:prior_loss}
\mathcal{L}_{\text{prior}}(\boldsymbol{\theta}) = \mathbb{E}_{\mathbf{z}_I, \mathbf{z}_c, t, \boldsymbol{\epsilon}} \left\| \boldsymbol{\epsilon} - \boldsymbol{\epsilon}_{\boldsymbol{\theta}}(\mathbf{z}_I^{(t)}, t, \mathbf{z}_c)\right\|_2^2.
\end{equation}
На инференсе априорная диффузионная модель генерирует оценку $\widehat{\mathbf{z}}_I \approx \mathbf{z}_I$ итеративной процедурой расшумления, согласующей $\mathbf{z}_c$ с визуальным пространством CLIP. Кодировщики $\mathcal{E}_f$, $\mathcal{E}_e$ и модуль слияния $\mathcal{F}$ остаются замороженными, что упрощает схему и стабилизирует обучение.

\paragraph{Стадия II: генерация изображения при условии $\widehat{\mathbf{z}}_I$.}
Для реконструкции изображения используется предобученная латентно-диффузионная модель SDXL~\cite{podell2023sdxlimprovinglatentdiffusion}:
\begin{equation*}
\boldsymbol{\epsilon}_{\boldsymbol{\phi}}(\mathbf{x}^{(t)}, t, h(\widehat{\mathbf{z}}_I)),
\end{equation*}
где $\mathbf{x}^{(t)}$~--- зашумленный латент изображения, а $h(\cdot)$~--- проекционный адаптер IP-Adapter~\cite{ye2023ipadaptertextcompatibleimage}, переводящий $\widehat{\mathbf{z}}_I$ в формат условия, совместимый с SDXL. Модель $\boldsymbol{\phi}$ не дообучается.

\paragraph{Итоговый процесс реконструкции.}
\begin{enumerate}
    \item Контрастивное обучение из раздела~\ref{sec3:contrastive} выравнивает $\mathbf{z}_c$ с $\mathbf{z}_I$.
    \item Априорная диффузионная модель генерирует согласованное визуальное представление $\widehat{\mathbf{z}}_I$ по $\mathbf{z}_c$.
    \item Предобученная SDXL, обусловленная через $h(\widehat{\mathbf{z}}_I)$, синтезирует изображение $\widehat{\mathbf{x}}$.
\end{enumerate}

\section{Вычислительный эксперимент}\label{sec4:bimodal}
В данном разделе проводится эмпирическая оценка метода, предложенного в разделе~\ref{sec3:bimodal}. Раздел~\ref{sec4:data} описывает характеристики используемого датасета, параметры регистрации сигналов и процедуры предобработки, включая восстановление отсутствующих каналов ЭЭГ и построение индивидуальных масок мозга для фМРТ. В разделе~\ref{sec4:hyper} приведены архитектурные детали и гиперпараметры энкодеров для каждой модальности, а также модуля объединения эмбеддингов. Раздел~\ref{sec4:results} содержит результаты вычислительных экспериментов: качественный анализ реконструкций и количественную оценку через метрику CLIP-Score. Ограничения предложенного метода обсуждаются в разделе~\ref{sec4:limits}. 
\subsection{Данные}\label{sec4:data}

В работе используется открытый датасет с одновременной регистрацией фМРТ и ЭЭГ при предъявлении визуальных стимулов~\cite{Telesford2023}. Данные собраны у 22 здоровых пациентов в возрасте 23--51 года в ходе двух сканирований. Эксперимент включает записи в состоянии покоя, при просмотре мерцающего шахматного поля, абстрактной анимации \texttt{Inscapes}~--- компьютерно сгенерированной последовательности трехмерных фигур с плавными переходами~--- а также коротких видеороликов, представляющих натуралистические стимулы: фрагменты мультфильмов \texttt{The Present} и \texttt{\texttt{Despicable Me}}, а также фрагменты фильмов с обезьянами. Все визуальные стимулы демонстрируются в естественном, натуралистическом формате, что способствует высокой межсубъектной корреляции и устойчивости нейронных откликов. Синхронизация ЭЭГ и фМРТ достигается через запись триггеров сканера в ЭЭГ-поток. Предобработка ЭЭГ включает удаление градиентных и кардиобаллистических артефактов, полосовую фильтрацию и усредненное реферирование. Для фМРТ применяется стандартная коррекция движения, сглаживание и нормализация. Параметры регистрации сведены в Таблице~\ref{table:params}.

Исходно данные поступают в виде трех асинхронных потоков: временные ряды фМРТ, временные ряды ЭЭГ и видеозапись стимулов. Регистрация ведется синхронно, однако частоты дискретизации у модальностей различны: фМРТ обновляется с частотой $0{,}476$~Гц, ЭЭГ записывается с частотой $5000$~Гц, а видеостимул образует непрерывный поток кадров. У каждого сигнала есть и собственные артефакты, у ЭЭГ градиентные и кардиобаллистические, у фМРТ двигательные, а набор доступных электродов у разных пациентов неодинаков; в совокупности это исключает прямую подачу сырых записей на батчевое обучение нейросетевой модели. Поэтому первый шаг предобработки~--- разбиение записей на согласованные триплеты $(\mathbf{F}, \mathbf{E}, \mathbf{I})$. С учетом гемодинамического лага $\Delta t = 5$~с, согласующегося с физиологическим диапазоном и результатом главы~\ref{ch:delay}, каждому фМРТ-тензору в момент времени $t$ ставятся в соответствие окно ЭЭГ и кадр видеостимула, отвечающие моменту $t - \Delta t$. Видеопоток предварительно разбивается на кадры с частотой 3 кадра в секунду (англ. Frames Per Second, FPS), а ЭЭГ-сигнал сегментируется на матрицы размера $(K, T)$, где $K$~--- число каналов, а $T$~--- длина временного окна.

Параллельно с временной синхронизацией проводится пространственная предобработка, важная для снижения размерности и стандартизации входных тензоров. Чтобы снизить размерность фМРТ и убрать фоновые воксели, для каждого пациента строится индивидуальная маска мозга~\eqref{eq:mask}: вычисляется суммарное абсолютное изменение сигнала во времени по каждому вокселю, и к полученному тензору применяется метод Оцу~\cite{otsu1975threshold} для автоматического выбора порога. Результат~--- бинарная маска мозга, отделяющая воксели мозга от фона. Пример применения маски мозга на срезы снимка фМРТ показан на Рисунке~\ref{fig:fMRI_masking}. 
Важный этап подготовки данных~--- приведение записей ЭЭГ к фиксированному набору каналов. Монтаж набора данных~\cite{Telesford2023} включает $K = 61$ кортикальный электрод. Расположение электродов на скальпе согласно системе 10--20 показано на Рисунке~\ref{fig:eeg_montage}. Однако в ряде записей часть каналов отсутствует, в результате чего количество доступных каналов $|\mathcal{C}_i| \ge 50$ варьируется. Пересечение множеств доступных каналов для всей выборки оказывается пустым ($\bigcap_i \mathcal{C}_i = \varnothing$), поскольку не существует канала, присутствующего без исключений во всех элементах выборки. Следовательно, использование общего для всех записей подмножества каналов невозможно, что обусловливает необходимость восстановления отсутствующих сигналов до полного набора из $K = 61$ канала согласно методу, описанному в разделе~\ref{sec3:recover}, см.~\eqref{eq:recover}.
Для выбора оптимальной стратегии восстановления проведено сравнение четырех методов: нулевого заполнения, интерполяции по ближайшему соседу и двух вариантов метода $k$-ближайших соседей с евклидовой и сферической метриками. Качество интерполяции оценивалось по среднеквадратичной ошибке (англ. Mean Squared Error, MSE) восстановления сигнала, рассчитанной независимо для каждого канала. Результаты представлены на Рисунке~\ref{fig:eeg_chan_recovery}, где по оси абсцисс отложен индекс канала, а по оси ординат~--- значение ошибки.

\begin{table}[h]
    \centering
    \caption{Описание оборудования и основные параметры нейросигналов}
    \label{table:params}
    \resizebox{\textwidth}{!}{\begin{tabular}{llll}
    \toprule
    Модальность & Оборудование & Частота дискретизации & Каналы \\
    \midrule
    ЭЭГ & Brain Products BrainCap MR & 5000~Гц & 61 кортикальный электрод \\
    фМРТ & Siemens 3.0~Т TIM Trio & 0,476~Гц & 12-канальная головная катушка \\
    \bottomrule
    \end{tabular}}
\end{table}

\begin{figure}[h!]
    \centering
    \includegraphics[width=1.\textwidth]{figure_4.png}
    \caption{Для снижения размерности данных фМРТ применяются маски головного мозга, исключающие фоновые воксели, не несущие информации о визуальном стимуле. Маски формируются на основе анализа наиболее изменчивых вокселей в ответ на стимуляцию. Воксели, включенные в маску, выделены красным цветом}
    \label{fig:fMRI_masking}
\end{figure}

\begin{figure}[h!]
    \centering
    \includegraphics[width=0.8\textwidth]{figs/eeg_montage.pdf}
    \caption{Конфигурация монтажа из $K = 61$ кортикального электрода, используемая в датасете~\cite{Telesford2023}. Датчики расположены согласно системе 10--20. Данная схема используется для восстановления отсутствующих каналов методом $k$-ближайших соседей, описанным в разделе~\ref{sec3:recover}.}
    \label{fig:eeg_montage}
\end{figure}

\begin{figure}[h!]
    \centering
    \includegraphics[width=1.\textwidth]{figs/eeg_recovery.pdf}
    \caption{Сравнение стратегий восстановления отсутствующих каналов ЭЭГ. По оси абсцисс отложен индекс канала, по оси ординат~--- среднеквадратичная ошибка (англ. Mean Squared Error, MSE). Метод нулевого заполнения характеризуется максимальной ошибкой, тогда как алгоритмы на основе $k$-ближайших соседей обеспечивают минимальные значения MSE для большинства каналов.}
    \label{fig:eeg_chan_recovery}
\end{figure}

\subsection{Параметры энкодеров}\label{sec4:hyper}

В экспериментах использовались две специализированные архитектуры, адаптированные под особенности каждой модальности. Энкодер фМРТ начинается с персонализированной линейной проекции: для каждого пациента воксели, отобранные с помощью индивидуальной бинарной маски, проецируются в пространство размерности 4096 обучаемым линейным слоем. За этим следует остаточный многослойный перцептрон, состоящий из трех полносвязных преобразований: первый расширяет представление до 2048, второй сохраняет эту размерность, а третий сжимает до 1024. Между слоями применяются нелинейность GELU (англ. Gaussian Error Linear Unit) и нормализация по слою; к выходу каждого блока добавляется остаточное соединение от входа. Доля прореживания (англ. dropout) во всех слоях равна 0{,}1.

Для обработки ЭЭГ-сигналов применялся компактный энкодер. Входные данные имеют размер $61 \times 525$, где 60 каналов соответствуют кортикальным электродам, а один дополнительный канал~--- обучаемому эмбеддингу пациента. Архитектура включает три последовательные свертки: временная свертка с ядром $(1, 32)$ и 256 фильтрами, пространственная свертка с ядром $(61, 1)$ и 512 фильтрами, а также глобальная свертка с ядром $(1, 525)$ и 1024 фильтрами. После сверток следует MLP-проектор, реализующий преобразование $1024 \to 2048 \to 1024$ с нелинейностью GELU и нормализацией по слою (англ. Layer Normalization, LayerNorm). Доля прореживания во всех сверточных и проекционных слоях равна 0{,}1. Выходной эмбеддинг имеет размерность 1024, что обеспечивает совместимость с модулем слияния и CLIP-пространством.

Модуль объединения эмбеддингов фМРТ и ЭЭГ реализован как двухслойный MLP: первый слой расширяет конкатенированное представление (размерность $2048$) до $4096$, второй сжимает его обратно до целевой размерности $1024$. Активация~--- GELU, доля прореживания~--- 0{,}1. Выходной эмбеддинг лежит в CLIP-пространстве и подается как условие при генерации изображений.

\subsection{Результаты экспериментов}\label{sec4:results}

Для вычислительных экспериментов из всего набора стимулов датасета~\cite{Telesford2023} было отобрано подмножество видеозаписей, разумно сочетающее объем данных с семантическим разнообразием визуальных стимулов. Всего датасет содержит пять типов видеостимулов: абстрактную анимацию \texttt{Inscapes}, мультфильмы \texttt{The Present} и \texttt{Despicable Me}, а также натуралистические видеоролики с обезьянами~--- \texttt{monkey1}, \texttt{monkey2}, \texttt{monkey5}. Два из них~--- \texttt{Inscapes} и \texttt{Despicable Me}~--- заметно сдвинуты по домену относительно реалистичных сцен. \texttt{Inscapes}~--- это компьютерно сгенерированная последовательность абстрактных трехмерных фигур, лишенных семантически интерпретируемых объектов, а \texttt{The Present} и \texttt{Despicable Me} принадлежат к жанру анимации с сильно стилизованной визуальной подачей. Обучение на гетерогенной смеси столь различных доменов без явной доменной адаптации приводит к противоречивым градиентам и ухудшению обобщающей способности модели. В связи с этим для основного эксперимента были выбраны три видеостимула из категории \texttt{monkey}, которые обеспечивают достаточное разнообразие сцен: различные ракурсы, освещение, действия персонажей, при сохранении единого визуального домена~--- натуралистической видеосъемки. Объемы данных по каждому датасету приведены в Таблице~\ref{tab:monkey_dataset}.

\begin{table}[h!]
    \centering
    \caption{Объемы данных по видеостимулам датасета \texttt{monkey*}, использованным в основном эксперименте. Данные разделены в пропорции 80\% к 20\% для обучения и тестирования соответственно, с равномерным распределением записей пациентов между выборками}
    \label{tab:monkey_dataset}
    \begin{tabular}{lccc}
        \toprule
        Видеостимул & Общий объем & Обучающая выборка & Тестовая выборка \\
        \midrule
        \texttt{monkey1} & 5\,941 &  4\,753 & 1\,188 \\
        \texttt{monkey2} & 4\,477 & 3\,582 & 895 \\
        \texttt{monkey5} & 4\,579 & 3\,664 & 915 \\
        \bottomrule
    \end{tabular}
\end{table}

Для верификации эффективности предложенного подхода проведены два эксперимента: визуальный анализ изображений-реконструкций и количественная оценка через CLIP-Score.

В первом эксперименте выполнено сравнение реконструкций, полученных с диффузионной априорной моделью и без нее на тестовых объектах из датасета \texttt{monkey*}. Примеры приведены на Рисунке~\ref{fig:reconstructions}. Для проверки гипотезы о чувствительности модели к распределению визуальных стимулов был проведен дополнительный эксперимент по кросс-доменной реконструкции. Модель, обученная исключительно на датасете \texttt{monkey*}, применялась для декодирования сигналов фМРТ--ЭЭГ, зарегистрированных при просмотре абстрактной анимации \texttt{Inscapes}. Поскольку визуальные характеристики стимулов \texttt{Inscapes} принципиально отличаются от реалистичных сцен, на которых проводилось обучение, энкодер фМРТ--ЭЭГ не получает на вход паттерны нейронной активности, соответствующие выученному распределению. В результате комбинированный эмбеддинг $\mathbf{z}_c$ формируется в области латентного пространства, удаленной от кластеров, ассоциированных с обучающими стимулами, что приводит к генерации семантически несогласованных изображений, см. Рисунок~\ref{fig:domain_shift}. Без явной доменной адаптации или аугментации тренировочных данных обобщение на распределения с сильным визуальным сдвигом не гарантируется. Для оценки устойчивости моделей необходим датасет с выраженной доменной разнородностью, однако, как показано в разделе~\ref{sec1:datasets}, открытые мультимодальные наборы данных, обладающие подобным разнообразием визуальных стимулов, на текущий момент отсутствуют.
\begin{figure}[h!]
    \centering
    \includegraphics[width=1.\textwidth]{figure_5.pdf}
    \caption{Примеры реконструкции изображений по сигналам фМРТ--ЭЭГ. Строки соответствуют различным режимам: (1)~оригинальное изображение; (2)~реконструкция из CLIP-эмбеддинга оригинального изображения с помощью предобученной SDXL как предельное качество реконструкции; (3)~реконструкция предложенным методом без диффузионной априорной модели; (4)~реконструкция предложенным методом с диффузионной априорной моделью. Применение априорной модели повышает семантическую согласованность реконструированных изображений с исходными стимулами}
    \label{fig:reconstructions}
\end{figure}
\begin{figure}[h!]
    \centering
    \includegraphics[width=0.45\textwidth]{figs/domain_shift_example.pdf}
    \caption{Пример реконструкции стимула из домена \texttt{Inscapes} моделью, обученной на датасете \texttt{monkey*}. Строки соответствуют различным режимам: (1)~оригинальное изображение; (2)~реконструкция из CLIP-эмбеддинга оригинального изображения с помощью предобученной SDXL как предельное качество реконструкции; (3)~реконструкция предложенным методом без диффузионной априорной модели; (4)~реконструкция предложенным методом с диффузионной априорной моделью. Отсутствие семантической согласованности обусловлено доменным сдвигом между обучающими и тестируемыми стимулами}
    \label{fig:domain_shift}
\end{figure}

Во втором эксперименте проводится количественная оценка с использованием метрики CLIP-Score~--- косинусной близости между CLIP-эмбеддингом реконструированного изображения $\mathbf{z}_I$ и комбинированным представлением мозговой активности $\mathbf{z}_c$:
\begin{equation}\label{eq:clip_score}
    \text{CLIP-Score}(\mathbf{z}_c, \mathbf{z}_I) = \frac{\langle \mathbf{z}_c, \mathbf{z}_I \rangle}{\| \mathbf{z}_c \| \cdot \| \mathbf{z}_I \|}.
\end{equation}
Метрики вычислялись на тестовой выборке датасета~\cite{Telesford2023}, использовались данные из видеостимулов \texttt{monkey1}, \texttt{monkey2}, \texttt{monkey5}. Усреднение проводилось по пациентам и объектам. Для получения результатов в Таблице~\ref{tab:clip_scores} были использованы следующие конфигурации моделей:
\begin{enumerate}
    \item \textbf{ЭЭГ}: используется только энкодер ЭЭГ $\mathcal{E}_e$, модуль объединения отключен, эмбеддинг $\mathbf{Z}_e$ напрямую используется для реконструкции без априорной модели.
    \item \textbf{фМРТ}: используется только энкодер фМРТ $\mathcal{E}_f$, модуль объединения отключен, эмбеддинг $\mathbf{Z}_f$ напрямую используется для реконструкции без априорной модели.
    \item \textbf{фМРТ--ЭЭГ}: используются оба энкодера $\mathcal{E}_e$ и $\mathcal{E}_f$ с модулем объединения $\mathcal{F}$, формирующим комбинированный эмбеддинг $\mathbf{Z}_c$ без априорной модели.
    \item \textbf{ЭЭГ \& Diffusion Prior}: используется энкодер ЭЭГ с последующей априорной диффузионной моделью для уточнения эмбеддинга перед генерацией.
    \item \textbf{фМРТ \& Diffusion Prior}: используется энкодер фМРТ с последующей априорной диффузионной моделью для уточнения эмбеддинга перед генерацией.
    \item \textbf{фМРТ--ЭЭГ \& Diffusion Prior}: используется полная архитектура с обоими энкодерами, модулем объединения и априорной диффузионной моделью.
\end{enumerate}

Полученные результаты демонстрируют, что использование одной модальности дает худшие показатели по сравнению с комбинированной моделью на основе фМРТ и ЭЭГ. Более того, интеграция диффузионной априорной модели существенно повышает CLIP-Score для всех конфигураций, а наилучшее качество достигается при совместном использовании фМРТ, ЭЭГ и диффузионной априорной модели. Это подтверждает, что мультимодальность компенсирует ограничения отдельных методов регистрации, а двухстадийная генерация необходима для семантической точности реконструкции.

\begin{table}[h!]
    \centering
    \caption{Сравнение CLIP-Score для различных конфигураций модели на трех видеостимулах из набора данных. Произведено усреднение по всем 22 пациентам, нижним регистром указаны среднеквадратичные отклонения. Наилучшие результаты выделены жирным шрифтом}
    \label{tab:clip_scores}
    \begin{tabular}{lccc}
        \toprule
        & \multicolumn{3}{c}{Название видеостимула} \\
        \cline{2-4}
        Модель & \texttt{monkey1} & \texttt{monkey2} & \texttt{monkey5} \\
        \midrule
        ЭЭГ & $0{,}231_{\pm 0{,}004}$ & $0{,}262_{\pm 0{,}004}$ & $0{,}224_{\pm 0{,}005}$ \\
        фМРТ & $0{,}254_{\pm 0{,}004}$ & $0{,}270_{\pm 0{,}003}$ & $0{,}242_{\pm 0{,}003}$ \\
        фМРТ--ЭЭГ & $0{,}310_{\pm 0{,}003}$ & $0{,}312_{\pm 0{,}004}$ & $0{,}284_{\pm 0{,}003}$ \\
        ЭЭГ \& Diffusion Prior & $0{,}542_{\pm 0{,}005}$ & $0{,}557_{\pm 0{,}003}$ & $0{,}538_{\pm 0{,}003}$ \\
        фМРТ \& Diffusion Prior & $0{,}564_{\pm 0{,}005}$ & $0{,}571_{\pm 0{,}004}$ & $0{,}551_{\pm 0{,}005}$ \\
        фМРТ--ЭЭГ \& Diffusion Prior & $\textbf{0{,}668}_{\pm 0{,}003}$ & $\textbf{0{,}647}_{\pm 0{,}005}$ & $\textbf{0{,}644}_{\pm 0{,}003}$ \\
        \bottomrule
    \end{tabular}
\end{table}

\subsection{Ограничения метода}\label{sec4:limits}

Отдельно стоит оговорить ограничения, которые накладывает использование CLIP-пространства как промежуточного представления. CLIP обучался сопоставлять изображения с текстовыми описаниями, поэтому его эмбеддинги несут преимущественно семантику высокого уровня~--- объекты, их взаимосвязи, общий контекст сцены~--- но почти не удерживают низкоуровневые визуальные характеристики: точные текстуры, локальные паттерны, пиксельные детали. Как следствие, предложенный метод реконструкции нацелен на семантическую согласованность полученных изображений с исходными стимулами, а не на точное пиксельное воспроизведение низкоуровневых деталей. Это ограничение вытекает из выбора CLIP в качестве промежуточного представления и присуще всем подходам на основе CLIP-эмбеддингов~\cite{huang2021fmri, li2024visual, scotti2023reconstructingmindseyefmritoimage}.

Метод проверялся на одном публичном датасете~\cite{Telesford2023}, а качество реконструкции оценивалось преимущественно метрикой CLIP-Score. Кросс-доменные эксперименты подтвердили, что обобщающая способность модели ограничена: на стимулах из домена, которого не было в обучающей выборке~--- например, на абстрактной анимации \texttt{Inscapes}~--- качество реконструкции заметно снижается. Энкодеры подстраивают свои представления под конкретное распределение визуальных признаков обучающего домена и не обладают встроенной инвариантностью к резким семантическим сдвигам. Поэтому перенос метода на другие типы стимулов~--- статичные изображения, абстрактные паттерны, текстовые сцены~--- либо на иные условия регистрации потребует дополнительной валидации и, вероятно, явной доменной адаптации. Сказывается и индивидуальная вариабельность нейронной активности пациентов, от которой также зависит качество метода. Привлечение гетерогенных датасетов или альтернативных метрик усилило бы выводы, хотя и имеющихся результатов достаточно для подтверждения работоспособности подхода в рамках заявленного домена.

% TODO \fixme{Дополнить главу: условия, при которых мультимодальное объединение строго улучшает одномодальное декодирование; эксперименты на мультимодальных временных рядах иной природы для демонстрации переносимости.}
